% ============================================================
\chapter{Modul 9: Objektorientierung}
% ============================================================

\section{Lektion 24: Warum Klassen?}

\begin{analogie}
Ein Roboter ist kein Bündel loser Funktionen -- er ist ein Objekt mit Eigenschaften (Akkustand, Position, Modus) und Fähigkeiten (fahren, drehen, stoppen). Eine Klasse beschreibt dieses Objekt. Ein konkreter Roboter ist eine Instanz dieser Klasse.
\end{analogie}

\section{Lektionen 25--28: Der komplette Roboter als Klasse}

\begin{lstlisting}
import math, time

class MobilRoboter:
    """Differentialantrieb-Roboter mit Ultraschallsensor."""

    RADABSTAND  = 0.15   # Meter
    MAX_V       = 0.50   # m/s

    def __init__(self, name, akku=100):
        self.name     = name
        self.akku     = akku
        self.x        = 0.0
        self.y        = 0.0
        self.theta    = 0.0   # Orientierung in Rad
        self.modus    = "bereit"
        self._laufzeit = 0.0

    # ── Bewegung ─────────────────────────────────────────────
    def fahre(self, v, omega, dt=0.1):
        """Bewegt den Roboter nach Kinematik-Gleichungen."""
        v     = max(-self.MAX_V, min(self.MAX_V, v))
        dx    = v * math.cos(self.theta) * dt
        dy    = v * math.sin(self.theta) * dt
        dtheta = omega * dt
        self.x     += dx
        self.y     += dy
        self.theta += dtheta
        self.akku  -= abs(v) * dt * 2
        self._laufzeit += dt

    def drehe_grad(self, winkel_deg, omega=0.5):
        """Dreht um einen bestimmten Winkel."""
        winkel_rad = math.radians(abs(winkel_deg))
        dt = winkel_rad / omega
        richtung = 1 if winkel_deg > 0 else -1
        self.fahre(0, richtung * omega, dt)

    # ── Sensorik (simuliert) ──────────────────────────────────
    def lese_abstand(self):
        import random
        return round(random.uniform(0.1, 3.0), 2)

    # ── Status ───────────────────────────────────────────────
    @property
    def position(self):
        return (round(self.x, 3), round(self.y, 3))

    @property
    def orientierung_grad(self):
        return round(math.degrees(self.theta), 1)

    def status(self):
        print(f"{'='*40}")
        print(f"  Roboter:     {self.name}")
        print(f"  Position:    {self.position}")
        print(f"  Orientierung:{self.orientierung_grad} Grad")
        print(f"  Akku:        {self.akku:.1f}%")
        print(f"  Modus:       {self.modus}")
        print(f"  Laufzeit:    {self._laufzeit:.1f} s")
        print(f"{'='*40}")

    def __repr__(self):
        return f"MobilRoboter('{self.name}', akku={self.akku:.0f}%)"


class FolgenderRoboter(MobilRoboter):
    """Erweiterung: folgt einem Zielroboter."""

    def __init__(self, name, folgeziel=None):
        super().__init__(name)
        self.folgeziel = folgeziel

    def folge(self, ziel_x, ziel_y, dt=0.1):
        dx = ziel_x - self.x
        dy = ziel_y - self.y
        d  = math.sqrt(dx**2 + dy**2)
        if d > 0.3:
            v     = min(0.3, d * 0.5)
            theta_ziel = math.atan2(dy, dx)
            omega = (theta_ziel - self.theta) * 2.0
            self.fahre(v, omega, dt)


# ── Demonstration ─────────────────────────────────────────────
if __name__ == "__main__":
    r1 = MobilRoboter("Alpha")
    r2 = FolgenderRoboter("Beta")

    print("Starte Fahrt...")
    for _ in range(20):
        abstand = r1.lese_abstand()
        if abstand > 0.25:
            r1.fahre(0.3, 0.0)
        else:
            r1.drehe_grad(90)

        r2.folge(r1.x, r1.y)

    r1.status()
    r2.status()
    print(repr(r1))
\end{lstlisting}

\begin{merksatz}
\texttt{@property} macht eine Methode wie ein Attribut nutzbar: \texttt{r.position} statt \texttt{r.position()}. Das ist der pythonische Weg für berechnete Eigenschaften. \texttt{\_variable} (ein Unterstrich) signalisiert: intern, bitte nicht direkt ändern.
\end{merksatz}

\begin{uebung}[title=Projektaufgabe Modul 9: Roboterflotte]
Erstelle eine Klasse \texttt{Sensor} mit Unterklassen \texttt{Ultraschall} und \texttt{IR\_Sensor}. Binde sie in \texttt{MobilRoboter} ein: \texttt{self.sensoren = \{"vorne": Ultraschall(), "links": IR\_Sensor(), ...\}}. Schreibe eine Methode \texttt{alle\_sensoren\_lesen()}, die über das Sensor-Dictionary iteriert.
\end{uebung}
